% ==========================================================
%  nhanxet.tex  --  100% English
% ==========================================================

% ------------------------------------------------------------------
% DECLARATION
% ------------------------------------------------------------------
\newpage
\chapter*{DECLARATION}
\addcontentsline{toc}{chapter}{DECLARATION}

I, \textbf{\StudentName} (Student ID: \textbf{\StudentID}, Class: \textbf{\ClassCode}), hereby solemnly declare that this thesis is my own original research work carried out under the guidance of \textbf{\SupervisorTitle{} \SupervisorName} at Can Tho University.

All results, data, figures, and tables presented in this thesis are truthful and have not been previously submitted for any academic degree or diploma at this or any other institution. All external sources, ideas, and findings that are not my own have been properly cited in accordance with academic standards.

\vspace{1.5cm}
\begin{flushright}
    Can Tho, \textit{Month Day}, 20\ldots\\[8pt]
    \textit{Author}\\[2.5cm]
    \textbf{\StudentName}
\end{flushright}

% ------------------------------------------------------------------
% ACKNOWLEDGEMENTS
% ------------------------------------------------------------------
\newpage
\chapter*{ACKNOWLEDGEMENTS}
\addcontentsline{toc}{chapter}{ACKNOWLEDGEMENTS}

First and foremost, I would like to express my deepest gratitude to my supervisor, \textbf{\SupervisorTitle{} \SupervisorName}, for the invaluable guidance, patience, and constructive feedback throughout the entire duration of this thesis. Despite a demanding schedule, my supervisor consistently offered insightful direction and encouragement that kept this project on track.

I am sincerely grateful to the faculty members of the Department of Software Engineering, College of Information and Communication Technology, Can Tho University, for providing an excellent academic environment and for generously sharing their expertise throughout my studies.

My heartfelt thanks go to the deaf and hard-of-hearing community members at Can Tho University who participated in requirement-gathering sessions and provided invaluable feedback that shaped the user interface and data-collection workflow of \textbf{CTU-SignBridge}.

Finally, I owe an immeasurable debt of gratitude to my family and close friends for their unwavering support, patience, and encouragement throughout my academic journey.

\vspace{1cm}
\begin{flushright}
    \textbf{\StudentName}
\end{flushright}

% ------------------------------------------------------------------
% ABSTRACT
% ------------------------------------------------------------------
\newpage
\chapter*{ABSTRACT}
\addcontentsline{toc}{chapter}{ABSTRACT}

Vietnamese Sign Language (VSL) exhibits significant variations arising from regional, cultural, and educational differences across Vietnam, posing substantial challenges for AI-based recognition systems that rely on standardized, homogeneous training datasets. Furthermore, no dedicated infrastructure currently exists to support the collaborative, multi-group collection of dialect-varied VSL samples at scale. Existing tools are either repurposed general-purpose platforms that lack structured annotation support, or informal file-sharing arrangements that produce data loss under network failure, undetected cross-group vocabulary duplication, and incompatible labeling conventions --- all of which critically impede reproducible research on multi-dialect VSL.

To address these limitations, this thesis presents \textbf{CTU-SignBridge}, a purpose-built, self-hosted platform designed to serve as the data infrastructure layer for multi-dialect VSL research. CTU-SignBridge contributes three tightly integrated subsystems:

First, a \textbf{multi-tenant data collection subsystem} is developed based on a Workspace--Project architecture with Role-Based Access Control (RBAC). This subsystem supports the acquisition of VSL samples from multiple dialect groups and individual contributors while enforcing complete data isolation between independent research teams on shared infrastructure. Duplicate submissions are eliminated through client-side SHA-256 checksum verification before any data reaches the server.

Second, a \textbf{structured vocabulary taxonomy subsystem} is designed using a Star Schema model (LANGUAGES $\to$ DIALECTS $\to$ CLASSES, extended by CATEGORIES and SIGN\_FEATURES satellite tables). This design explicitly models the multi-dialect structure of VSL, supports faceted multi-criteria filtering without recursive tree traversal, and prevents retroactive label corruption through a Taxonomy Forking Engine.

Third, a \textbf{landmark-first data ingestion pipeline} is implemented using Google MediaPipe Holistic running entirely within the browser. The pipeline extracts 543 skeletal keypoints per frame in real time and compresses them into \texttt{.npz} files, achieving a 99.1\% storage reduction compared to raw MP4 video while preserving all gesture information necessary for downstream recognition model training.

% ---- RESULTS PLACEHOLDER: replace with actual numbers after experiments ----
\textit{[Experimental results: to be completed after evaluation.
Planned metrics: API P95 response latency at 500 concurrent users;
end-to-end upload latency per session;
storage reduction ratio (.npz vs.\ raw MP4);
Celery retry success rate under simulated worker failure;
multi-tenant isolation verification.]}
% ---------------------------------------------------------------------------

CTU-SignBridge provides the data-collection foundation required for community-driven, scalable, and reproducible Vietnamese Sign Language research, reducing the barrier to large-scale VSL dataset construction on commodity hardware.

\vspace{0.5cm}
\noindent\textbf{Keywords:} Vietnamese Sign Language, Data Collection Platform, Multi-tenancy, Star Schema, MediaPipe Holistic, FastAPI, Celery, Redis, RBAC.
